\documentclass[11pt,a4paper]{article}
\usepackage{listings}
\lstdefinestyle{framed}{basicstyle=\ttfamily\small, columns=fullflexible,
  numbers=left, numberstyle=\tiny, frame=single}
\begin{document}
\section{Code}
Listing~\ref{lst:py} is set with the package's defaults, in fixed columns,
and so is the inline \lstinline|x = f(y)| in this paragraph.
\begin{lstlisting}[language=Python, caption={Squares in Python.}, label=lst:py]
def squares(n):
    # the first n squares
    return [i * i for i in range(n)]
\end{lstlisting}
The text continues after the first listing, long enough to reach a second
line of this paragraph.

\begin{lstlisting}[language=C, style=framed]
int main(void) {
    printf("hello, world\n");
    return 0;
}
\end{lstlisting}

After both listings the text continues.
\end{document}
